% Options for packages loaded elsewhere
\PassOptionsToPackage{unicode}{hyperref}
\PassOptionsToPackage{hyphens}{url}
\documentclass[
  11pt,
]{article}
\usepackage{xcolor}
\usepackage[margin=1in]{geometry}
\usepackage{amsmath,amssymb}
\setcounter{secnumdepth}{-\maxdimen} % remove section numbering
\usepackage{iftex}
\ifPDFTeX
  \usepackage[T1]{fontenc}
  \usepackage[utf8]{inputenc}
  \usepackage{textcomp} % provide euro and other symbols
\else % if luatex or xetex
  \usepackage{unicode-math} % this also loads fontspec
  \defaultfontfeatures{Scale=MatchLowercase}
  \defaultfontfeatures[\rmfamily]{Ligatures=TeX,Scale=1}
\fi
\usepackage{lmodern}
\ifPDFTeX\else
  % xetex/luatex font selection
\fi
% Use upquote if available, for straight quotes in verbatim environments
\IfFileExists{upquote.sty}{\usepackage{upquote}}{}
\IfFileExists{microtype.sty}{% use microtype if available
  \usepackage[]{microtype}
  \UseMicrotypeSet[protrusion]{basicmath} % disable protrusion for tt fonts
}{}
\makeatletter
\@ifundefined{KOMAClassName}{% if non-KOMA class
  \IfFileExists{parskip.sty}{%
    \usepackage{parskip}
  }{% else
    \setlength{\parindent}{0pt}
    \setlength{\parskip}{6pt plus 2pt minus 1pt}}
}{% if KOMA class
  \KOMAoptions{parskip=half}}
\makeatother
\setlength{\emergencystretch}{3em} % prevent overfull lines
\providecommand{\tightlist}{%
  \setlength{\itemsep}{0pt}\setlength{\parskip}{0pt}}
\usepackage{mathtools}
\usepackage{amsthm}
\usepackage{mathrsfs}
\usepackage{enumitem}
\usepackage{microtype}
\usepackage{fancyhdr}
\usepackage{titlesec}

\setlist{nosep}
\setcounter{secnumdepth}{3}
\setcounter{tocdepth}{2}
\setlength{\headheight}{14pt}
\emergencystretch=4em

\pagestyle{fancy}
\fancyhf{}
\lhead{Short-Kappa Branch}
\rhead{\thepage}
\renewcommand{\headrulewidth}{0.4pt}

\titleformat{\section}
  {\normalfont\Large\bfseries}
  {\thesection}
  {0.75em}
  {}
\titleformat{\subsection}
  {\normalfont\large\bfseries}
  {\thesubsection}
  {0.75em}
  {}

\newcommand{\rethlaslabel}[1]{\textnormal{\small\texttt{#1}}}
\usepackage{bookmark}
\IfFileExists{xurl.sty}{\usepackage{xurl}}{} % add URL line breaks if available
\urlstyle{same}
\hypersetup{
  pdftitle={Short-Kappa Branch of Maximal-Parabolic Induction},
  pdfauthor={Rethlas verified blueprint},
  hidelinks,
  pdfcreator={LaTeX via pandoc}}

\title{Short-Kappa Branch of Maximal-Parabolic Induction}
\author{Rethlas verified blueprint}
\date{May 4, 2026}

\begin{document}
\maketitle

\section{lemma lem:x0-geometry-and-k}\label{lemma-lemx0-geometry-and-k}

\subsection{statement}\label{statement}

One has \[
(\operatorname{Im} X_0)^\perp=\ker X_0.
\] Hence the restriction of \(\langle\cdot,\cdot\rangle_0\) to
\(\ker X_0\) has radical \[
R:=\ker X_0\cap \operatorname{Im} X_0,
\] and therefore induces a non-degenerate \(\epsilon\)-symmetric form on
\[
K_{X_0}:=\ker X_0/R.
\]

Moreover there exists a subspace \[
V_1\subset \ker X_0
\] such that the quotient map \[
q:\ker X_0\to K_{X_0}
\] restricts to an isometric isomorphism \[
q|_{V_1}:V_1\xrightarrow{\sim} K_{X_0}.
\] For such a choice one has an orthogonal decomposition \[
\ker X_0=V_1\oplus R
\] and \[
V_0=V_1\oplus V_1^\perp,
\] where \(V_1^\perp\) is \(X_0\)-stable.

Finally, if
\(h\in \operatorname{Isom}(K_{X_0},\langle\cdot,\cdot\rangle_K)\), then
there is an element \[
\widetilde h\in Z_{G_0}(X_0)
\] whose restriction to \(V_1\) is \((q|_{V_1})^{-1} h (q|_{V_1})\) and
which acts trivially on \(V_1^\perp\).

\subsection{proof}\label{proof}

If \(w\in \ker X_0\) and \(y=X_0 z\in \operatorname{Im} X_0\), then \[
\langle w,y\rangle_0=\langle w,X_0 z\rangle_0=-\langle X_0 w,z\rangle_0=0.
\] So \[
\ker X_0\subseteq (\operatorname{Im} X_0)^\perp.
\] Since \[
\dim (\operatorname{Im} X_0)^\perp
=\dim V_0-\dim \operatorname{Im} X_0
=\dim \ker X_0,
\] the inclusion is an equality.

Therefore the radical of the restriction of
\(\langle\cdot,\cdot\rangle_0\) to \(\ker X_0\) is \[
\ker X_0\cap (\ker X_0)^\perp
=\ker X_0\cap \operatorname{Im} X_0
=R,
\] so the form descends to a non-degenerate \(\epsilon\)-symmetric form
on \(K_{X_0}\).

Choose any linear section \(s:K_{X_0}\to \ker X_0\) of \(q\). The
pulled-back form \[
(a,b)\longmapsto \langle s(a),s(b)\rangle_0
\] is exactly \(\langle a,b\rangle_K\), hence is non-degenerate.
Therefore the image \[
V_1:=s(K_{X_0})
\] is a non-degenerate subspace of \(\ker X_0\), \(q|_{V_1}\) is an
isometric isomorphism, and \[
\ker X_0=V_1\oplus R
\] orthogonally.

Because \(V_1\) is non-degenerate in \(V_0\), we have \[
V_0=V_1\oplus V_1^\perp.
\] If \(v\in V_1^\perp\) and \(u\in V_1\), then \(X_0u=0\), so \[
\langle X_0v,u\rangle_0=-\langle v,X_0u\rangle_0=0.
\] Thus \(X_0v\in V_1^\perp\), proving that \(V_1^\perp\) is
\(X_0\)-stable.

Now let \(h\in \operatorname{Isom}(K_{X_0})\). Via the isometry
\(q|_{V_1}:V_1\to K_{X_0}\), it defines an isometry of \(V_1\), still
denoted \(h\). Extend it to \(V_0\) by letting it act trivially on
\(V_1^\perp\). Since \(V_0=V_1\oplus V_1^\perp\) is orthogonal, this
extension is an element \(\widetilde h\in G_0\).

Because \(X_0\) vanishes on \(V_1\) and preserves \(V_1^\perp\), and
because \(\widetilde h\) acts as the identity on \(V_1^\perp\), we have
\[
\widetilde h X_0 = X_0 \widetilde h.
\] So \(\widetilde h\in Z_{G_0}(X_0)\), as required.

\section{lemma
lem:parametrize-x0-plus-u}\label{lemma-lemparametrize-x0-plus-u}

\subsection{statement}\label{statement-1}

For \(C\in \operatorname{Hom}_F(E',V_0)\) let \(C^\vee:V_0\to E\) be
defined by \[
\lambda(C^\vee v,e')=-\langle v,Ce'\rangle_0.
\] For \(B\in \operatorname{Hom}_F(E',E)\) satisfying \[
\lambda(Bu,v)+\epsilon \lambda(Bv,u)=0
\qquad (u,v\in E'),
\] define \[
X_{C,B}(e+v+e'):=C^\vee v+Be'+X_0 v+Ce'.
\] Then \(X_{C,B}\in X_0+\mathfrak u\).

Conversely, every element of \(X_0+\mathfrak u\) is uniquely of the form
\(X_{C,B}\).

\subsection{proof}\label{proof-1}

Relative to the decomposition \(V=E\oplus V_0\oplus E'\), an element
\(Y\in \mathfrak u\) has the shape \[
Y=
\begin{pmatrix}
0&A&B\\
0&0&C\\
0&0&0
\end{pmatrix}
\] with \(A:V_0\to E\), \(B:E'\to E\), \(C:E'\to V_0\). The condition
\(Y\in \mathfrak g\) is \[
\langle Yx,y\rangle+\langle x,Yy\rangle=0.
\] Testing this on \((v,e')\in V_0\times E'\) gives \[
\lambda(Av,e')+\langle v,Ce'\rangle_0=0,
\] so \(A=C^\vee\). Testing it on \((u,v)\in E'\times E'\) gives \[
\lambda(Bu,v)+\epsilon\lambda(Bv,u)=0.
\] Hence every \(Y\in \mathfrak u\) is uniquely of the form \[
Y_{C,B}(e+v+e')=C^\vee v+Be'+Ce'.
\] Therefore every element of \(X_0+\mathfrak u\) is uniquely \[
X_0+Y_{C,B}=X_{C,B}.
\]

\section{lemma
lem:unipotent-conjugation}\label{lemma-lemunipotent-conjugation}

\subsection{statement}\label{statement-2}

For \(D\in \operatorname{Hom}_F(E',V_0)\) define \[
u_D(e+v+e')
:=\left(e+D^\vee v+\frac12 D^\vee D\,e'\right)+(v+De')+e'.
\] Then \(u_D\in P\). Moreover, for every pair \((C,B)\) as above, \[
u_D X_{C,B} u_D^{-1}
=X_{C-X_0D,\;B+D^\vee C-C^\vee D-D^\vee X_0D}.
\]

\subsection{proof}\label{proof-2}

In block form, \[
u_D=
\begin{pmatrix}
1&D^\vee&\frac12 D^\vee D\\
0&1&D\\
0&0&1
\end{pmatrix},
\] and its inverse is \[
u_D^{-1}=
\begin{pmatrix}
1&-D^\vee&\frac12 D^\vee D\\
0&1&-D\\
0&0&1
\end{pmatrix}.
\] To check that \(u_D\) preserves the form, it is enough to test it on
pairs involving the \(E'\)-summand.

For \(v\in V_0\) and \(e'\in E'\), \[
\langle D^\vee v,e'\rangle=\lambda(D^\vee v,e')=-\langle v,De'\rangle_0,
\] so the mixed \((V_0,E')\)-terms cancel.

For \(u',v'\in E'\), \begin{align*}
\langle u_Du',u_Dv'\rangle
&=\lambda\!\left(\frac12 D^\vee D\,u',v'\right)
+\epsilon\lambda\!\left(\frac12 D^\vee D\,v',u'\right)
+\langle Du',Dv'\rangle_0 \\
&=-\frac12\langle Du',Dv'\rangle_0
-\frac{\epsilon}{2}\langle Dv',Du'\rangle_0
+\langle Du',Dv'\rangle_0 \\
&=0,
\end{align*} because
\(\langle Dv',Du'\rangle_0=\epsilon\langle Du',Dv'\rangle_0\). Hence
\(u_D\in G\), and since it fixes \(E\) and stabilizes \(E^\perp\), in
fact \(u_D\in P\).

Now \[
X_{C,B}=
\begin{pmatrix}
0&C^\vee&B\\
0&X_0&C\\
0&0&0
\end{pmatrix}.
\] Multiplying the three block matrices \(u_D\), \(X_{C,B}\), and
\(u_D^{-1}\) gives \begin{align*}
u_DX_{C,B}u_D^{-1}
&=
\begin{pmatrix}
0&C^\vee+D^\vee X_0&B+D^\vee C-C^\vee D-D^\vee X_0D\\
0&X_0&C-X_0D\\
0&0&0
\end{pmatrix}.
\end{align*} Since \[
\bigl(C-X_0D\bigr)^\vee=C^\vee+D^\vee X_0
\] by the defining relation for \((-)^\vee\), this is exactly \[
X_{C-X_0D,\;B+D^\vee C-C^\vee D-D^\vee X_0D}.
\] This proves the formula.

\section{proposition
prop:generic-normal-form}\label{proposition-propgeneric-normal-form}

\subsection{statement}\label{statement-3}

Fix a \(d\)-dimensional subspace \[
A\subset V_1\simeq K_{X_0}.
\] Let \[
W_A\subset V_0/\operatorname{Im} X_0
\] be the annihilator of \(A\) under the perfect pairing \[
(V_0/\operatorname{Im} X_0)\times \ker X_0\to F,
\qquad (\overline v,k)\mapsto \langle k,v\rangle_0.
\] Then \(\dim W_A=\kappa\).

Choose any \(\kappa\)-dimensional subspace \(V_A\subset V_0\) mapping
isomorphically to \(W_A\), and any isomorphism \[
S_A:E'\xrightarrow{\sim} V_A.
\] Set \[
x_A:=X_{S_A,0}\in X_0+\mathfrak u.
\]

Then:

\begin{enumerate}
\def\labelenumi{\arabic{enumi}.}
\item
  The restriction \[
  S_A^\vee|_{\ker X_0}:\ker X_0\to E
  \] is surjective, with kernel exactly \(A\).
\item
  Let \(x=X_{C,B}\in X_0+\mathfrak u\), and suppose that \[
  \phi_C:=C^\vee|_{\ker X_0}:\ker X_0\to E
  \] is surjective. Put \[
  L:=\ker \phi_C.
  \] Assume that \[
  L\cap R=0.
  \] Then \(q|_L:L\xrightarrow{\sim} q(L)\subset K_{X_0}\) is an
  isomorphism onto a \(d\)-dimensional subspace \(A:=q(L)\), and \(x\)
  is \(P\)-conjugate to \(x_A\).
\end{enumerate}

Moreover, for a fixed \(A\), the \(P\)-orbit of \(x_A\) is independent
of the auxiliary choices of \(V_A\) and \(S_A\).

\begin{enumerate}
\def\labelenumi{\arabic{enumi}.}
\setcounter{enumi}{2}
\tightlist
\item
  If \(A_1,A_2\subset V_1\simeq K_{X_0}\) are isometric inside
  \(K_{X_0}\), then \(x_{A_1}\) and \(x_{A_2}\) are \(P\)-conjugate,
  hence certainly \(G\)-conjugate.
\end{enumerate}

\subsection{proof}\label{proof-3}

Since \(A\subset \ker X_0\) has dimension \(d\), its annihilator in
\(V_0/\operatorname{Im} X_0\) has dimension \[
c_1-d=\kappa,
\] which proves \(\dim W_A=\kappa\).

Now let \(a\in A\). By definition of \(W_A\), one has \[
\langle a,S_Ae'\rangle_0=0\qquad \forall e'\in E',
\] so \(S_A^\vee a=0\). Thus \[
A\subseteq \ker(S_A^\vee|_{\ker X_0}).
\] Conversely, if \(k\in \ker X_0\) satisfies \(S_A^\vee k=0\), then
\(k\) pairs trivially with \(V_A\). Because \[
\ker X_0=(\operatorname{Im} X_0)^\perp,
\] the value \(\langle k,v\rangle_0\) depends only on the class of \(v\)
in \(V_0/\operatorname{Im} X_0\). Since \(V_A\) maps isomorphically onto
\(W_A\), the condition \(S_A^\vee k=0\) says exactly that \(k\)
annihilates \(W_A\) under the perfect pairing \[
(V_0/\operatorname{Im} X_0)\times \ker X_0\to F.
\] But \(W_A\) was defined as the annihilator of \(A\) in this perfect
pairing, so the double-annihilator property gives \[
k\in \operatorname{Ann}(W_A)=A.
\] Thus the kernel is exactly \(A\). Since \[
\dim \ker X_0-\dim A=c_1-d=\kappa=\dim E,
\] the map \(S_A^\vee|_{\ker X_0}\) is surjective. This proves (1).

For (2), surjectivity of \(\phi_C\) gives \[
\dim L=c_1-\kappa=d.
\] The assumption \(L\cap R=0\) implies that \(q|_L\) is injective, so
\[
A:=q(L)
\] is a \(d\)-dimensional subspace of \(K_{X_0}\).

We first normalize the restriction of \(C^\vee\) to \(\ker X_0\). Both
maps \[
\phi_C:\ker X_0\to E
\qquad\text{and}\qquad
S_A^\vee|_{\ker X_0}:\ker X_0\to E
\] are surjective and have the same kernel \(L\). Hence there exists a
unique \(a\in \operatorname{GL}(E)\) such that \[
a\circ \phi_C = S_A^\vee|_{\ker X_0}.
\] Let \(d_E\in \operatorname{GL}(E')\) be the unique element satisfying
\[
a=(d_E^{-1})^\vee.
\] Conjugating by the corresponding Levi element of \(P\), we may
replace \((C,B)\) by a new pair for which already \[
C^\vee|_{\ker X_0}=S_A^\vee|_{\ker X_0}.
\]

Then \((C-S_A)^\vee\) vanishes on \(\ker X_0\). Equivalently, \[
\langle k,(C-S_A)e'\rangle_0=0
\qquad \forall k\in \ker X_0,\ \forall e'\in E'.
\] Since \((\ker X_0)^\perp=\operatorname{Im} X_0\), it follows that \[
(C-S_A)(E')\subseteq \operatorname{Im} X_0.
\] Hence there exists \(D_1:E'\to V_0\) with \[
X_0D_1=C-S_A.
\] By Lemma \texttt{lem:unipotent-conjugation}, conjugation by
\(u_{D_1}\) replaces \(C\) by \(S_A\).

So we may assume from now on that \(x=X_{S_A,B}\). Because
\(S_A^\vee|_{\ker X_0}\) is surjective, choose \(D_2:E'\to \ker X_0\)
such that \[
S_A^\vee D_2=\frac12 B.
\] Since \(D_2(E')\subset \ker X_0\), one has \(X_0D_2=0\), so
conjugation by \(u_{D_2}\) keeps \(C=S_A\) fixed. The new \(B\)-block is
\[
B' = B + D_2^\vee S_A - S_A^\vee D_2.
\] For \(u,v\in E'\), \[
\lambda(D_2^\vee S_Au,v)=-\langle S_Au,D_2v\rangle_0.
\] Because \(S_A^\vee D_2=\frac12 B\), \[
\langle D_2u,S_Av\rangle_0
=-\lambda\!\left(\frac12 Bu,v\right).
\] Using
\(\langle S_Au,D_2v\rangle_0=\epsilon\langle D_2v,S_Au\rangle_0\) and \[
\lambda(Bv,u)=-\epsilon\lambda(Bu,v),
\] we get \[
\lambda(D_2^\vee S_Au,v)=-\frac12\lambda(Bu,v).
\] Also \[
\lambda(S_A^\vee D_2u,v)=\frac12\lambda(Bu,v).
\] Therefore \[
\lambda(B'u,v)=0
\qquad \forall u,v\in E',
\] so \(B'=0\). Thus \(x\) is \(P\)-conjugate to \(x_A\).

Now fix \(A\) and let \(\widetilde V_A,\widetilde S_A\) be another
choice with image in \(W_A\), and put \[
\widetilde x_A:=X_{\widetilde S_A,0}.
\] By (1), the map \[
\widetilde S_A^\vee|_{\ker X_0}
\] is surjective with kernel \(A\). Since \[
A\subset V_1
\qquad\text{and}\qquad
V_1\cap R=0,
\] we have \[
A\cap R=0.
\] Therefore Part (2), applied to \(\widetilde x_A\), shows that \[
\widetilde x_A
\] is \(P\)-conjugate to the already chosen \(x_A\). This proves the
choice-independence statement.

For (3), let \(h:A_1\to A_2\) be induced by an element of
\(\operatorname{Isom}(K_{X_0})\). By Lemma
\texttt{lem:x0-geometry-and-k}, \(h\) lifts to an element
\(\widetilde h\in Z_{G_0}(X_0)\). Extend \(\widetilde h\) to a Levi
element of \(P\) as follows. Let \[
W_{A_2}':=\widetilde h(W_{A_1})\subset V_0/\operatorname{Im} X_0,
\] choose \[
V_{A_2}':=\widetilde h(V_{A_1})\subset V_0,
\] and define \[
S_{A_2}':=\widetilde h\circ S_{A_1}:E'\xrightarrow{\sim} V_{A_2}'.
\] Then \[
V_{A_2}'
\] maps isomorphically onto \(W_{A_2}\), and the Levi element built from
\(\widetilde h\) on \(V_0\), from the induced map on \(E'\), and from
the dual inverse on \(E\), sends \[
x_{A_1}
\] exactly to \[
X_{S_{A_2}',0}.
\] By the choice-independence already proved, this latter element is
\(P\)-conjugate to the originally chosen representative \(x_{A_2}\).
Hence \(x_{A_1}\) and \(x_{A_2}\) are \(P\)-conjugate.

\section{lemma
lem:maximal-rank-classes}\label{lemma-lemmaximal-rank-classes}

\subsection{statement}\label{statement-4}

Let \(\mathscr A^{\mathrm{mr}}\) be the maximal-rank locus inside the
Grassmannian of \(d\)-dimensional subspaces of the real non-degenerate
\(\epsilon\)-space \(K_{X_0}\).

Then the possible isometry classes of elements of
\(\mathscr A^{\mathrm{mr}}\) are finite, and each isometry class is open
in \(\mathscr A^{\mathrm{mr}}\).

More precisely:

\begin{enumerate}
\def\labelenumi{\arabic{enumi}.}
\item
  If \(\epsilon=-1\) and \(d\) is even, then every element of
  \(\mathscr A^{\mathrm{mr}}\) is a non-degenerate symplectic space of
  dimension \(d\), hence the isometry class is unique.
\item
  If \(\epsilon=-1\) and \(d\) is odd, then every element of
  \(\mathscr A^{\mathrm{mr}}\) is a symplectic space of rank \(d-1\)
  with one-dimensional radical, hence the isometry class is unique.
\item
  If \(\epsilon=+1\), then every element of \(\mathscr A^{\mathrm{mr}}\)
  is a non-degenerate symmetric real bilinear space of dimension \(d\),
  hence is classified by its signature. In particular only finitely many
  signatures occur, and each signature class is open.
\end{enumerate}

\subsection{proof}\label{proof-4}

If \(\epsilon=-1\), the restricted form is alternating. On an
even-dimensional space a non-degenerate alternating form is unique up to
isometry, and on an odd-dimensional space an alternating form of maximal
rank \(d-1\) is unique up to isometry: in a suitable basis it is
standard symplectic on a \((d-1)\)-dimensional subspace and zero on a
radical line.

We also need ambient transitivity in these two symplectic cases. If
\(\epsilon=-1\) and \(d\) is even, let
\(A,B\in \mathscr A^{\mathrm{mr}}\). Choose symplectic bases \[
a_1,b_1,\dots,a_{d/2},b_{d/2}
\] of \(A\) and \[
a_1',b_1',\dots,a_{d/2}',b_{d/2}'
\] of \(B\), and extend both to symplectic bases of the ambient space
\(K_{X_0}\). The basis-to-basis map is an element of \[
\operatorname{Isom}(K_{X_0},\langle\cdot,\cdot\rangle_K)
\] carrying \(A\) to \(B\).

If \(\epsilon=-1\) and \(d\) is odd, let \(r\) and \(r'\) span the
radicals of \(A\) and \(B\). Choose symplectic bases \[
a_1,b_1,\dots,a_{(d-1)/2},b_{(d-1)/2}
\] of a complement to \(Fr\) in \(A\), and \[
a_1',b_1',\dots,a_{(d-1)/2}',b_{(d-1)/2}'
\] of a complement to \(Fr'\) in \(B\). Choose vectors
\(s,s'\in K_{X_0}\) such that \[
\langle r,s\rangle_K=1,
\qquad
\langle r',s'\rangle_K=1,
\] and extend \[
r,a_1,b_1,\dots,a_{(d-1)/2},b_{(d-1)/2},s
\] and \[
r',a_1',b_1',\dots,a_{(d-1)/2}',b_{(d-1)/2}',s'
\] to symplectic bases of \(K_{X_0}\). Again the basis-to-basis map is
an ambient isometry carrying \(A\) to \(B\). Thus the ambient isometry
orbit is unique in both symplectic cases.

If \(\epsilon=+1\), the restricted form is symmetric. Maximal rank then
means non-degenerate. Over \(\mathbb R\), such a form is classified by
its signature, and the signature is locally constant on the
non-degenerate locus. So only finitely many signatures occur, each open
in \(\mathscr A^{\mathrm{mr}}\).

It remains to see that, for a fixed signature, the corresponding
\(d\)-planes form a single orbit under
\(\operatorname{Isom}(K_{X_0},\langle\cdot,\cdot\rangle_K)\). Let
\(A,B\subset K_{X_0}\) be non-degenerate \(d\)-dimensional subspaces
with the same signature. Then there is an isometry \[
u:A\xrightarrow{\sim} B.
\] Because \(A\) and \(B\) are non-degenerate, the orthogonal
complements \[
A^\perp,\qquad B^\perp
\] are also non-degenerate, and their signatures are the signature of
\(K_{X_0}\) minus the common signature of \(A\) and \(B\). Hence there
is an isometry \[
v:A^\perp\xrightarrow{\sim} B^\perp.
\] Taking orthogonal direct sums gives an ambient isometry \[
u\oplus v\in \operatorname{Isom}(K_{X_0},\langle\cdot,\cdot\rangle_K)
\] with \[
(u\oplus v)(A)=B.
\] Thus every fixed-signature locus is a single ambient isometry orbit.

\section{proposition
prop:kernel-image-quotient}\label{proposition-propkernel-image-quotient}

\subsection{statement}\label{statement-5}

For \(A\subset V_1\simeq K_{X_0}\) of dimension \(d\), let \(x_A\) be as
in Proposition \texttt{prop:generic-normal-form}. Then \[
\ker x_A=E\oplus A,
\] \[
\operatorname{Im} x_A=E\oplus \operatorname{Im} X_0\oplus S_A(E'),
\] and \[
\ker x_A\cap \operatorname{Im} x_A = E\oplus \operatorname{rad}(A).
\]

Consequently, \[
\ker x_A/(\ker x_A\cap \operatorname{Im} x_A)\simeq A/\operatorname{rad}(A),
\] and the induced form on this quotient is exactly the quotient of
\(\langle\cdot,\cdot\rangle_K|_A\) by its radical.

In particular, if \(A\in \mathscr A^{\mathrm{mr}}\) and either
\(\epsilon=+1\) or \(d\) is even, then \(A\) is non-degenerate and \[
\ker x_A\cap \operatorname{Im} x_A = E,
\qquad
\ker x_A/(\ker x_A\cap \operatorname{Im} x_A)\simeq A.
\]

\subsection{proof}\label{proof-5}

Take \[
y=e+v+e'\in E\oplus V_0\oplus E'.
\] Then \[
x_A(y)=S_A^\vee v+X_0v+S_Ae'.
\] If \(x_A(y)=0\), then the \(V_0\)-component gives \[
X_0v+S_Ae'=0.
\] The class of \(S_Ae'\) in \(V_0/\operatorname{Im} X_0\) is zero only
if \(e'=0\), so \(v\in \ker X_0\). Then the \(E\)-component gives \[
S_A^\vee v=0,
\] hence \(v\in A\) by Proposition \texttt{prop:generic-normal-form}(1).
So \[
\ker x_A=E\oplus A.
\]

The formula for the image is immediate from \[
x_A(e+v+e')=S_A^\vee v+X_0v+S_Ae':
\] the \(E\)-part is \(S_A^\vee(\ker X_0)=E\), the
\(\operatorname{Im} X_0\)-part comes from \(X_0(V_0)\), and the
remaining quotient part comes from \(S_A(E')\).

Now let \[
z=e+a\in E\oplus A=\ker x_A.
\] Then \(z\in \operatorname{Im} x_A\) if and only if \[
a\in A\cap (\operatorname{Im} X_0+S_A(E')).
\] We claim that, for \(a\in A\), this is equivalent to the class \[
\overline a\in V_0/\operatorname{Im} X_0
\] belonging to \(W_A\). Indeed, if \[
a=X_0v+S_Ae',
\] then \(\overline a=\overline{S_Ae'}\in W_A\). Conversely, if
\(\overline a\in W_A\), choose \(e'\in E'\) with \[
\overline{S_Ae'}=\overline a;
\] then \(a-S_Ae'\in \operatorname{Im} X_0\), so \[
a\in \operatorname{Im} X_0+S_A(E').
\]

Now suppose first that \[
a\in A\cap (\operatorname{Im} X_0+S_A(E')).
\] Then \(\overline a\in W_A\). For every \(b\in A\), the definition of
\(W_A\) gives \[
0=\langle b,a\rangle_0.
\] Since \(a,b\in V_1\subset \ker X_0\), this is the same as \[
0=\langle b,a\rangle_K.
\] Hence \(a\in A^\perp\cap A=\operatorname{rad}(A)\).

Conversely, let \(a\in \operatorname{rad}(A)\). Then for every
\(b\in A\), \[
\langle b,a\rangle_K=0,
\] so again \[
\langle b,a\rangle_0=0.
\] Thus the class \(\overline a\) annihilates \(A\), hence \[
\overline a\in W_A.
\] By the claim above this is equivalent to \[
a\in \operatorname{Im} X_0+S_A(E').
\] Therefore \[
\ker x_A\cap \operatorname{Im} x_A = E\oplus \operatorname{rad}(A).
\]

Since \(E\) is orthogonal to \(V_0\), the form on \(\ker x_A=E\oplus A\)
induces exactly the quotient form of \(\langle\cdot,\cdot\rangle_K|_A\)
on \[
(E\oplus A)/(E\oplus \operatorname{rad}(A))
\simeq A/\operatorname{rad}(A).
\] The final statement is immediate when \(A\) is non-degenerate.

\section{proposition
prop:induced-orbits}\label{proposition-propinduced-orbits}

\subsection{statement}\label{statement-6}

The maximal \(G\)-orbits in \[
\operatorname{Ind}_P^G(\mathcal O_0)
\] are exactly the \(G\)-orbits of the representatives \(x_A\) with \[
A\in \mathscr A^{\mathrm{mr}},
\] modulo isometry inside \((K_{X_0},\langle\cdot,\cdot\rangle_K)\).

Equivalently, if \(A_1,\dots,A_N\) are representatives for the isometry
classes in \(\mathscr A^{\mathrm{mr}}\), then the maximal induced
\(G\)-orbits are exactly \[
\mathcal O_i:=G\cdot x_{A_i}
\qquad (1\le i\le N),
\] and these orbits are pairwise distinct.

\subsection{proof}\label{proof-6}

Embed \(G_0\) into the Levi factor of \(P\) by \[
\iota(g_0)(e+v+e'):=e+g_0v+e'.
\] Then \[
\iota(g_0)X_0\iota(g_0)^{-1}=g_0X_0g_0^{-1}.
\] Also \(\mathfrak u\) is \(P\)-stable: if \(p\in P\), then \(p(E)=E\)
and \(p(E^\perp)=E^\perp\), so the conditions \[
Y(E)=0,\qquad Y(E^\perp)\subseteq E,\qquad Y(V)\subseteq E^\perp
\] are preserved by \(\operatorname{Ad}(p)\). Therefore \[
\mathcal O_0+\mathfrak u
=\bigcup_{g_0\in G_0}\bigl(\iota(g_0)X_0\iota(g_0)^{-1}+\mathfrak u\bigr)
=\iota(G_0)\cdot (X_0+\mathfrak u),
\] and hence \[
G\cdot(\mathcal O_0+\mathfrak u)=G\cdot(X_0+\mathfrak u).
\] So \[
\operatorname{Ind}_P^G(\mathcal O_0)=\overline{G\cdot(X_0+\mathfrak u)}.
\]

We next isolate an open dense subset of the slice. For \[
x=X_{C,B}\in X_0+\mathfrak u
\] define \[
\phi_C:=C^\vee|_{\ker X_0}:\ker X_0\to E.
\] Let \(\Omega\subset X_0+\mathfrak u\) be the subset of all
\(X_{C,B}\) such that:

\begin{enumerate}
\def\labelenumi{\arabic{enumi}.}
\tightlist
\item
  \(\phi_C\) is surjective;
\item
  \(L:=\ker \phi_C\) satisfies \(L\cap R=0\);
\item
  the image \(A:=q(L)\subset K_{X_0}\) lies in
  \(\mathscr A^{\mathrm{mr}}\).
\end{enumerate}

This set is open and dense.

Indeed, under the parametrization of Lemma
\texttt{lem:parametrize-x0-plus-u}, the slice is the product of the
affine space \[
\operatorname{Hom}(E',V_0)
\] with the affine space of admissible \(B\)-blocks, and the defining
conditions for \(\Omega\) depend only on \(C\).

The assignment \[
\rho:\operatorname{Hom}(E',V_0)\longrightarrow \operatorname{Hom}(\ker X_0,E),
\qquad
\rho(C)=\phi_C,
\] is linear and surjective. Indeed, for \(k\in \ker X_0\) one has \[
\lambda(\phi_C(k),e')=-\langle k,Ce'\rangle_0,
\] so \(\phi_C\) depends only on the class of \(C\) in \[
\operatorname{Hom}(E',V_0/\operatorname{Im} X_0).
\] Using the perfect pairing \[
(V_0/\operatorname{Im} X_0)\times \ker X_0\to F,
\] one gets \[
V_0/\operatorname{Im} X_0\simeq (\ker X_0)^\vee,
\] and then the pairing \(\lambda:E\times E'\to F\) identifies \[
\operatorname{Hom}(E',V_0/\operatorname{Im} X_0)
\simeq
\operatorname{Hom}(\ker X_0,E).
\] Thus \(\rho\) is the quotient map to \[
\operatorname{Hom}(E',V_0/\operatorname{Im} X_0)
\] followed by that isomorphism.

Let \[
\operatorname{Surj}:=
\{\phi\in \operatorname{Hom}(\ker X_0,E):\phi\ \text{surjective}\}.
\] This is the usual open dense full-rank locus. Put \[
\operatorname{Surj}_R:=
\{\phi\in \operatorname{Surj}:\phi|_R\ \text{injective}\}.
\] Since \[
\ker \phi\cap R=0
\iff
\phi|_R\ \text{injective},
\] condition (2) is exactly the requirement that
\(\phi_C\in \operatorname{Surj}_R\). Because
\(c_2=\dim R\le \kappa=\dim E\), injectivity on \(R\) is again a
maximal-rank condition, hence open and dense. Therefore \[
\rho^{-1}(\operatorname{Surj}_R)\times
\{\text{admissible }B\}
\] is an open dense subset of \(X_0+\mathfrak u\).

Now define \[
\alpha:\operatorname{Surj}_R\to \operatorname{Gr}_d(K_{X_0}),
\qquad
\alpha(\phi)=q(\ker \phi).
\] This is well defined because for \(\phi\in \operatorname{Surj}_R\)
one has \[
\dim \ker \phi=c_1-\kappa=d
\] and \[
\ker \phi\cap R=0.
\]

We claim that \(\alpha\) is a surjective locally trivial map with affine
fibers. First fix \[
\phi_0\in \operatorname{Surj}_R,
\qquad
L_0:=\ker \phi_0,
\] and choose a complement \(M_0\) to \(L_0\) in \(\ker X_0\). On the
open neighborhood \[
\mathcal U_0:=
\{\phi\in \operatorname{Surj}_R:\phi|_{M_0}\ \text{is an isomorphism}\}
\] every kernel is the graph of the unique linear map \[
T_\phi:=-(\phi|_{M_0})^{-1}\circ \phi|_{L_0}:L_0\to M_0.
\] Thus \[
\phi\longmapsto \ker \phi
\] is locally trivial on \(\operatorname{Surj}_R\).

Next use \[
\ker X_0=V_1\oplus R.
\] If \(L\subset \ker X_0\) is \(d\)-dimensional and satisfies
\(L\cap R=0\), then the projection \[
\operatorname{pr}_{V_1}:L\to V_1
\] is injective, hence identifies \(L\) with a unique \(d\)-plane \[
A_L:=\operatorname{pr}_{V_1}(L)\subset V_1
\] and a unique linear map \[
\tau_L:A_L\to R
\] such that \[
L=\{a+\tau_L(a):a\in A_L\}.
\] Under the isometric identification \[
q|_{V_1}:V_1\xrightarrow{\sim} K_{X_0},
\] one has \[
q(L)=q(A_L)=A_L.
\] Therefore the map \[
L\longmapsto q(L)
\] from the open subset \[
\{L\in \operatorname{Gr}_d(\ker X_0):L\cap R=0\}
\] to \(\operatorname{Gr}_d(K_{X_0})\simeq \operatorname{Gr}_d(V_1)\)
is, in the usual graph charts, a trivial affine bundle with fiber
\(\operatorname{Hom}(A,R)\). Composing this affine bundle with the
locally trivial kernel map proves the claim about \(\alpha\).

The map \(\alpha\) is surjective: given any \[
A\in \operatorname{Gr}_d(K_{X_0}),
\] view it inside \(V_1\) via \(q|_{V_1}^{-1}\) and choose a complement
\(M\) to \(A\) in \(V_1\). Since \[
\dim(M\oplus R)=\kappa=\dim E,
\] there exists an isomorphism \[
\psi:M\oplus R\xrightarrow{\sim} E.
\] Extending \(\psi\) by \(0\) on \(A\) gives a surjective map \[
\phi:\ker X_0=A\oplus M\oplus R\to E
\] with \[
\ker \phi=A
\qquad\text{and}\qquad
\phi|_R\ \text{injective}.
\] Hence \(\phi\in \operatorname{Surj}_R\) and \(\alpha(\phi)=A\).

Because \(\mathscr A^{\mathrm{mr}}\) is an open dense subset of \[
\operatorname{Gr}_d(K_{X_0}),
\] its pullback \[
\alpha^{-1}(\mathscr A^{\mathrm{mr}})
\] is open dense in \(\operatorname{Surj}_R\). Therefore \[
\Omega
=
\{X_{C,B}:\phi_C\in \alpha^{-1}(\mathscr A^{\mathrm{mr}})\}
\] is open dense in \(X_0+\mathfrak u\).

On this locus the normalized slice datum may indeed be read, after the
\(P\)-conjugations of Proposition \texttt{prop:generic-normal-form},
from a map \[
S:E'\to \ker X_0,
\] not necessarily with image in \(V_1\), and the natural generic
condition is that the Gram matrix of \(\langle\cdot,\cdot\rangle_0\) on
\[
\operatorname{im}(S)
\] have maximal possible rank. The final orbit parameter, however, is
not \[
\operatorname{im}(S),
\] but the quotient subspace \[
A=q(\ker \phi_C)\subset K_{X_0},
\] and condition (3) is exactly that intrinsic maximal-rank condition.

By Proposition \texttt{prop:generic-normal-form}, every point of
\(\Omega\) is \(P\)-conjugate to some \(x_A\) with
\(A\in \mathscr A^{\mathrm{mr}}\). By Proposition
\texttt{prop:generic-normal-form}(3), isometric subspaces give the same
\(P\)-orbit. By Lemma \texttt{lem:maximal-rank-classes}, only finitely
many isometry classes occur, and each class is open in
\(\mathscr A^{\mathrm{mr}}\). Choose representatives \(A_1,\dots,A_N\)
for these classes, and let \[
\Omega_i:=P\cdot x_{A_i}.
\] Then \[
\Omega=\bigsqcup_{i=1}^N \Omega_i,
\] and each \(\Omega_i\) is open in \(\Omega\).

Set \[
\mathcal O_i:=G\cdot x_{A_i}.
\] Then \[
G\cdot \Omega = \bigcup_{i=1}^N \mathcal O_i.
\] Fix \(i\) and write \[
x_i:=x_{A_i}.
\] Because \[
\Omega_i=P\cdot x_i
\] is open in the affine slice \(X_0+\mathfrak u\), its tangent space at
\(x_i\) is \[
T_{x_i}(X_0+\mathfrak u)=T_{x_i}(P\cdot x_i)=[\mathfrak p,x_i]=\mathfrak u.
\] Now consider the map \[
f:G\times (X_0+\mathfrak u)\to \mathfrak g,
\qquad
f(g,z)=gzg^{-1}.
\] Its differential at \((1,x_i)\) is \[
df_{(1,x_i)}(\xi,v)=[\xi,x_i]+v
\qquad
(\xi\in \mathfrak g,\ v\in \mathfrak u).
\] Since \[
v\in \mathfrak u=[\mathfrak p,x_i]\subseteq [\mathfrak g,x_i],
\] the image of this differential is exactly \[
[\mathfrak g,x_i]=T_{x_i}(\mathcal O_i).
\] Therefore \(G\cdot (X_0+\mathfrak u)\) has, near \(x_i\), the same
local dimension as the orbit \(\mathcal O_i\), and the latter is open in
\(G\cdot (X_0+\mathfrak u)\) near \(x_i\). Conjugating by elements of
\(G\) shows that \(\mathcal O_i\) is open at each of its points, hence
open in \[
G\cdot (X_0+\mathfrak u).
\]

Since \(\Omega\) is dense in \(X_0+\mathfrak u\), the subset
\(G\times \Omega\) is dense in \(G\times (X_0+\mathfrak u)\), so \[
G\cdot \Omega
\] is dense in \(G\cdot (X_0+\mathfrak u)\). Therefore each
\(\mathcal O_i\) is an orbit of maximal dimension inside
\(G\cdot (X_0+\mathfrak u)\).

Because \[
\operatorname{Ind}_P^G(\mathcal O_0)
=\overline{G\cdot(X_0+\mathfrak u)}
=\overline{G\cdot \Omega},
\] and because \[
G\cdot \Omega=\bigcup_{i=1}^N \mathcal O_i
\] is a finite union, we get \[
\operatorname{Ind}_P^G(\mathcal O_0)
=\bigcup_{i=1}^N \overline{\mathcal O_i}.
\] Thus if \(\mathcal O\) is any \(G\)-orbit in the induced set and \[
x\in \mathcal O,
\] then \[
x\in \overline{\mathcal O_i}
\] for some \(i\). Since \(\overline{\mathcal O_i}\) is \(G\)-stable,
the whole orbit \(\mathcal O=G\cdot x\) is contained in
\(\overline{\mathcal O_i}\). Therefore every orbit in the induced set is
below at least one of the listed \(\mathcal O_i\) in the closure order.
It remains to prove that each listed orbit is itself maximal.

Set \[
Y:=G\cdot(X_0+\mathfrak u).
\] For fixed \(i\), the slice piece \[
\Omega_i=P\cdot x_i
\] is exactly the set of all \(x=X_{C,B}\in X_0+\mathfrak u\) such that:

\begin{enumerate}
\def\labelenumi{\arabic{enumi}.}
\tightlist
\item
  \(\phi_C=C^\vee|_{\ker X_0}\) is surjective;
\item
  \(L=\ker \phi_C\) satisfies \(L\cap R=0\);
\item
  the associated subspace \[
  A(x):=q(L)\subset K_{X_0}
  \] has maximal possible Gram-matrix rank and lies in the same real
  isometry class as \(A_i\).
\end{enumerate}

Indeed, conditions (1) and (2) are the defining open conditions for the
normalized slice, condition (3) is the intrinsic maximal-rank condition
on the kernel quotient parameter \(A=q(L)\), and Proposition
\texttt{prop:generic-normal-form} together with Lemma
\texttt{lem:maximal-rank-classes} says precisely that the resulting
points form the single \(P\)-orbit \(P\cdot x_i\). Since real isometry
class is locally constant on the maximal-rank locus, each \(\Omega_i\)
is open in \(X_0+\mathfrak u\). By the differential argument above, the
corresponding \(G\)-orbit \[
\mathcal O_i=G\cdot \Omega_i
\] is therefore open in \(Y\).

Now let \(\mathcal O\) be a \(G\)-orbit in \[
\operatorname{Ind}_P^G(\mathcal O_0)=\overline Y
\] such that \[
\mathcal O_i\subseteq \overline{\mathcal O}.
\] Since \(\mathcal O_i\) is open in \(Y\), it has the same dimension as
\(Y\): \[
\dim \mathcal O_i=\dim Y.
\] Because \(Y\) is dense in \[
\operatorname{Ind}_P^G(\mathcal O_0)=\overline Y,
\] the ambient induced set has the same dimension: \[
\dim \operatorname{Ind}_P^G(\mathcal O_0)=\dim Y.
\] Every \(G\)-orbit in \(\mathfrak g\) is a smooth locally closed
submanifold, so any orbit contained in the induced set has dimension at
most \[
\dim \operatorname{Ind}_P^G(\mathcal O_0)=\dim Y.
\] Hence \[
\dim \mathcal O\le \dim Y=\dim \mathcal O_i.
\] On the other hand, \[
\mathcal O_i\subseteq \overline{\mathcal O}
\] forces \[
\dim \mathcal O_i\le \dim \mathcal O.
\] Indeed, \(\mathcal O\) is locally closed, hence open in its closure
\(\overline{\mathcal O}\); therefore every orbit contained in the
boundary \[
\overline{\mathcal O}\setminus \mathcal O
\] has strictly smaller dimension than \(\mathcal O\).

Combining the two inequalities gives \[
\dim \mathcal O=\dim \mathcal O_i.
\] If \(\mathcal O\neq \mathcal O_i\), then \[
\mathcal O_i\subseteq \overline{\mathcal O}\setminus \mathcal O,
\] so the preceding boundary-dimension statement would imply \[
\dim \mathcal O_i<\dim \mathcal O,
\] contrary to the equality above. Therefore \[
\mathcal O=\mathcal O_i.
\] So each listed orbit \(\mathcal O_i\) is maximal.

It remains only to show that the \(\mathcal O_i\) are pairwise distinct
exactly according to the isometry class of \(A_i\). Suppose \[
G\cdot x_A = G\cdot x_{A'}
\] with \(A,A'\in \mathscr A^{\mathrm{mr}}\). The quotient \[
Q_x:=\ker x/(\ker x\cap \operatorname{Im} x)
\] is a \(G\)-orbit invariant bilinear space. By Proposition
\texttt{prop:kernel-image-quotient}, \[
Q_{x_A}\simeq A/\operatorname{rad}(A),
\qquad
Q_{x_{A'}}\simeq A'/\operatorname{rad}(A').
\]

If \(\epsilon=+1\), then every element of \(\mathscr A^{\mathrm{mr}}\)
is non-degenerate, so \(Q_{x_A}\simeq A\) and \(Q_{x_{A'}}\simeq A'\).
Hence \(A\) and \(A'\) are isometric.

If \(\epsilon=-1\) and \(d\) is even, the same conclusion holds because
maximal rank again means non-degenerate.

If \(\epsilon=-1\) and \(d\) is odd, Lemma
\texttt{lem:maximal-rank-classes} already says that all elements of
\(\mathscr A^{\mathrm{mr}}\) are isometric. So again the orbit depends
only on the isometry class of \(A\).

Therefore the maximal induced \(G\)-orbits are parametrized exactly by
the isometry classes in \(\mathscr A^{\mathrm{mr}}\).

\section{proposition
prop:multiplicity}\label{proposition-propmultiplicity}

\subsection{statement}\label{statement-7}

For every maximal induced orbit \(\mathcal O=G\cdot x_A\) with
\(A\in \mathscr A^{\mathrm{mr}}\), one has \[
m(\mathcal O,P)=
\begin{cases}
1,& \epsilon=+1,\\[0.4ex]
1,& \epsilon=-1\ \text{and}\ d\ \text{even},\\[0.4ex]
2,& \epsilon=-1\ \text{and}\ d\ \text{odd}.
\end{cases}
\]

\subsection{proof}\label{proof-7}

We treat first the case where \(A\) is non-degenerate. Then by
Proposition \texttt{prop:kernel-image-quotient}, \[
\ker x_A=E\oplus A,
\qquad
\ker x_A\cap \operatorname{Im} x_A=E.
\] Because \(A\) is non-degenerate and \(E\) is orthogonal to all of
\(\ker x_A\), the radical of the bilinear space \(\ker x_A\) is exactly
\(E\).

Any element of \(Z_G(x_A)\) preserves \(\ker x_A\), hence preserves its
radical. Therefore every element of \(Z_G(x_A)\) preserves \(E\), so \[
Z_G(x_A)\subseteq P.
\] Consequently \[
m(G\cdot x_A,P)=1.
\]

It remains to consider the case \(\epsilon=-1\) and \(d\) odd. Then
every \(A\in \mathscr A^{\mathrm{mr}}\) has one-dimensional radical.
Choose a symplectic decomposition \[
A=A_0\oplus Fr,
\] where \(A_0\) is non-degenerate symplectic and
\(Fr=\operatorname{rad}(A)\). We now construct a convenient
representative for this same \(A\) and then invoke the
choice-independence statement from Proposition
\texttt{prop:generic-normal-form}.

Choose \[
V_1=A_0\oplus Fr\oplus Fs\oplus B_0
\] with \[
\langle r,s\rangle_K=1
\] and \(B_0\) symplectic. Choose decompositions \[
E=Fe\oplus E_0,
\qquad
E'=Fe'\oplus E_0',
\] with \[
\lambda(e,e')=1
\] and \(E_0,E_0'\) paired perfectly. Let \[
W_0\subset V_0/\operatorname{Im} X_0
\] be the annihilator of \[
A\oplus Fs.
\] Then \[
\dim W_0=\kappa-1.
\] Choose a lift \[
V_0'\subset V_0
\] of \(W_0\) and an isomorphism \[
S_0:E_0'\xrightarrow{\sim} V_0'.
\] Let \(\overline r\) denote the class of \(r\) in
\(V_0/\operatorname{Im} X_0\). Since \(r\in \operatorname{rad}(A)\), one
has \[
\langle a,r\rangle_K=0
\qquad \forall a\in A,
\] so \[
\overline r\in W_A.
\] On the other hand, \[
\langle r,s\rangle_K=1,
\] so \(\overline r\) does not annihilate \(Fs\) and therefore \[
\overline r\notin W_0.
\] Because \[
\dim W_A=\kappa
\qquad\text{and}\qquad
\dim W_0=\kappa-1,
\] it follows that \[
W_A=W_0\oplus F\overline r.
\] Hence the subspace \[
\widetilde V_A:=Fr\oplus V_0'
\] maps isomorphically onto \(W_A\).

Define \[
\widetilde S_A:E'\xrightarrow{\sim} \widetilde V_A
\] by \[
\widetilde S_A(e')=r,
\qquad
\widetilde S_A|_{E_0'}=S_0.
\] and put \[
\widetilde x_A:=X_{\widetilde S_A,0}.
\] This is again a representative attached to the same subspace \(A\) in
Proposition \texttt{prop:generic-normal-form}. Therefore
\(\widetilde x_A\) is \(P\)-conjugate to the original \(x_A\), and \[
m(G\cdot x_A,P)=m(G\cdot \widetilde x_A,P).
\]

Set \[
U:=Fe\oplus Fr\oplus Fs\oplus Fe'.
\] Then \(U\) is non-degenerate symplectic,
\(\widetilde x_A(U)\subseteq U\), and in the ordered basis \[
e,\ r,\ s,\ e'
\] one has \[
\widetilde x_A(e)=0,\qquad \widetilde x_A(r)=0,\qquad \widetilde x_A(s)=e,\qquad \widetilde x_A(e')=r.
\] Because \(U\) is non-degenerate and \(\widetilde x_A\) is
skew-adjoint, the orthogonal complement \[
W:=U^\perp
\] is \(\widetilde x_A\)-stable. On \(U^\perp\) the kernel quotient is
\(A_0\), which is non-degenerate symplectic. Let \[
y:=\widetilde x_A|_W.
\] The first part of the present proof therefore applies to \(y\): every
element of its centralizer preserves the radical of \(\ker y\), namely
\[
E_0=E\cap W.
\] Thus \[
Z_{G(W)}(y)\subseteq P(W):=\{h\in G(W):h(E_0)=E_0\}.
\]

We now compute the residual \(4\times 4\) model on \(U\) correctly. Put
\[
U_0:=Fe\oplus Fr,
\qquad
U_1:=Fs\oplus Fe',
\qquad
U=U_0\oplus U_1.
\] Relative to the ordered basis \[
e,\ r,\ s,\ e',
\] the operator \(\widetilde x_A|_U\) is \[
\widetilde x_A|_U=
\begin{pmatrix}
0&I_2\\
0&0
\end{pmatrix},
\] and the symplectic form on \(U\) is represented by \[
J_U=
\begin{pmatrix}
0&S_2\\
-S_2&0
\end{pmatrix},
\qquad
S_2:=
\begin{pmatrix}
0&1\\
1&0
\end{pmatrix}.
\]

Let \[
G(U)=\operatorname{Sp}(U),
\qquad
P(U):=\{g\in G(U):g(Fe)=Fe\}.
\] Write an element of \(G(U)\) in \(2\times 2\) block form relative to
\[
U=U_0\oplus U_1
\] as \[
g=
\begin{pmatrix}
A&B\\
C&D
\end{pmatrix}.
\] The commutation relation \[
g\widetilde x_A|_U=\widetilde x_A|_Ug
\] is equivalent to \[
C=0,
\qquad
D=A.
\] So every element of \(Z_{G(U)}(\widetilde x_A|_U)\) has the form \[
g=
\begin{pmatrix}
A&B\\
0&A
\end{pmatrix}.
\] Now impose the symplectic condition \[
g^T J_U g=J_U.
\] Since \[
g=\begin{pmatrix}A&B\\0&A\end{pmatrix},
\qquad
J_U=\begin{pmatrix}0&S_2\\-S_2&0\end{pmatrix},
\] we first compute \[
J_U g
=
\begin{pmatrix}0&S_2\\-S_2&0\end{pmatrix}
\begin{pmatrix}A&B\\0&A\end{pmatrix}
=
\begin{pmatrix}0&S_2A\\-S_2A&-S_2B\end{pmatrix}.
\] Hence \[
g^T J_U g
=
\begin{pmatrix}A^T&0\\B^T&A^T\end{pmatrix}
\begin{pmatrix}0&S_2A\\-S_2A&-S_2B\end{pmatrix}
=
\begin{pmatrix}
0&A^T S_2A\\
-A^T S_2A&B^T S_2A-A^T S_2B
\end{pmatrix}.
\] Equating this with \[
J_U=\begin{pmatrix}0&S_2\\-S_2&0\end{pmatrix}
\] gives \[
A^T S_2 A=S_2
\] and \[
B^T S_2A=A^T S_2B.
\] If we write \[
M:=A^{-1}B,
\] then the second relation becomes \[
(AM)^T S_2A=A^T S_2(AM),
\] that is, \[
M^T A^T S_2A=A^T S_2A\,M.
\] Using \(A^T S_2A=S_2\), this is equivalent to \[
M^T S_2=S_2M.
\] Thus \(A\) lies in the orthogonal group \[
O(S_2):=\{A\in \operatorname{GL}_2(\mathbb R):A^T S_2 A=S_2\},
\] and for each fixed \(A\) the possible \(B\) form the affine space \[
A\cdot \{M\in M_2(\mathbb R):M^T S_2=S_2M\}.
\] Therefore projection to the diagonal block induces a surjective map
\[
Z_{G(U)}(\widetilde x_A|_U)\to O(S_2)
\] with affine fibers that are connected in the usual real topology, and
a section \[
A\longmapsto
\begin{pmatrix}
A&0\\
0&A
\end{pmatrix}.
\] Hence \[
\pi_0\bigl(Z_{G(U)}(\widetilde x_A|_U)\bigr)\simeq \pi_0(O(S_2)).
\]

Over \(\mathbb R\), the group \(O(S_2)\) admits the explicit
parametrization \[
O(S_2)=
\left\{
\begin{pmatrix}
a&0\\
0&a^{-1}
\end{pmatrix}:a\in \mathbb R^\times
\right\}
\sqcup
\left\{
\begin{pmatrix}
0&a\\
a^{-1}&0
\end{pmatrix}:a\in \mathbb R^\times
\right\}.
\] Thus there are two diagonal branches and two off-diagonal branches.
Each of the two diagonal branches lies in a different connected
component according to the sign of \(a\), and the two off-diagonal
branches also lies in a different connected component according to the
sign of \(a\). Hence there are four connected components in all, so \[
|\pi_0(O(S_2))|=4,
\] and hence \[
|\pi_0(Z_{G(U)}(\widetilde x_A|_U))|=4.
\]

If, in addition, \(g\) lies in \(P(U)\), then \(g\) preserves the line
\(Fe\subset U_0\). Because \(g\) acts on \(U_0\) by the block \(A\),
this means that \(A\) itself preserves \(Fe\). An element of \(O(S_2)\)
preserves \(Fe\) if and only if it has the form \[
\begin{pmatrix}
a&0\\
0&a^{-1}
\end{pmatrix}
\qquad (a\in \mathbb R^\times).
\] So projection to \(A\) induces a surjective map \[
Z_{G(U)}(\widetilde x_A|_U)\cap P(U)\to \mathbb R^\times
\] with the same real-topologically connected affine fibers as above and
with a section. Consequently \[
|\pi_0\bigl(Z_{G(U)}(\widetilde x_A|_U)\cap P(U)\bigr)|=2.
\] Indeed, inside the \(Fe\)-stabilizer only the diagonal branches
occur, with the two connected components given by \[
a>0
\qquad\text{and}\qquad
a<0.
\]

Finally, write an element of \(Z_G(\widetilde x_A)\), or of
\(Z_G(\widetilde x_A)\cap P\), in block form relative to the orthogonal
decomposition \[
V=U\oplus W.
\] Since \[
\widetilde x_A=\widetilde x_A|_U\oplus y,
\] write \[
g=
\begin{pmatrix}
a&b\\
c&d
\end{pmatrix},
\] where \[
a:U\to U,\qquad b:W\to U,\qquad c:U\to W,\qquad d:W\to W.
\] Projection to the diagonal blocks gives maps \[
Z_G(\widetilde x_A)\to Z_{G(U)}(\widetilde x_A|_U)\times Z_{G(W)}(y)
\] and \[
Z_G(\widetilde x_A)\cap P
\to
\bigl(Z_{G(U)}(\widetilde x_A|_U)\cap P(U)\bigr)\times Z_{G(W)}(y)
\] which are surjective because every pair in either base defines a
block-diagonal centralizing element. The second base has the same
\(W\)-factor as the first because \[
Z_{G(W)}(y)\subseteq P(W).
\] Indeed, relative to \[
E=Fe\oplus E_0
\] with \[
Fe\subset U,
\qquad
E_0\subset W,
\] a block-diagonal element \[
\operatorname{diag}(a,d)
\] lies in \(P\) exactly when \[
a\in P(U)
\qquad\text{and}\qquad
d\in P(W).
\] Since every \[
d\in Z_{G(W)}(y)
\] already lies in \(P(W)\), intersecting with \(P\) imposes no extra
condition on the \(W\)-factor and only cuts the \(U\)-factor down to \[
Z_{G(U)}(\widetilde x_A|_U)\cap P(U).
\]

Fix \((a,d)\) in the base of the first map. The fiber over \((a,d)\)
consists of all pairs \((b,c)\) such that \[
\begin{pmatrix}
a&b\\
c&d
\end{pmatrix}
\] lies in \(Z_G(\widetilde x_A)\). The commutation relation with
\(\widetilde x_A\) is equivalent to the linear equations \[
by=\widetilde x_A|_U\,b,
\qquad
c\widetilde x_A|_U=yc.
\] Let \((-)^*\) denote the adjoint relative to the forms on \(U\) and
\(W\). Because \(U\) and \(W\) are orthogonal, the condition \(g\in G\)
is equivalent to \[
a^*a+c^*c=1_U,
\qquad
a^*b+c^*d=0,
\qquad
b^*b+d^*d=1_W.
\] Since \(a\in G(U)\) and \(d\in G(W)\), this reduces to \[
c^*c=0,
\qquad
a^*b+c^*d=0,
\qquad
b^*b=0.
\] Therefore, if \((b,c)\) lies in the fiber, then for every
\(t\in[0,1]\) the pair \((tb,tc)\) satisfies the same equations. So each
fiber is star-shaped with center \((0,0)\), hence connected. The same
argument applies to the fibers of the second map.

Because both projections are surjective and admit the block-diagonal
section, they induce bijections on connected components. Thus \[
\pi_0(Z_G(\widetilde x_A))
\simeq
\pi_0\!\bigl(Z_{G(U)}(\widetilde x_A|_U)\times Z_{G(W)}(y)\bigr),
\] \[
\pi_0(Z_G(\widetilde x_A)\cap P)
\simeq
\pi_0\!\bigl((Z_{G(U)}(\widetilde x_A|_U)\cap P(U))\times Z_{G(W)}(y)\bigr).
\] Hence \[
|\pi_0(Z_G(\widetilde x_A))|
=
|\pi_0(Z_{G(U)}(\widetilde x_A|_U))|\cdot |\pi_0(Z_{G(W)}(y))|,
\] \[
|\pi_0(Z_G(\widetilde x_A)\cap P)|
=
|\pi_0(Z_{G(U)}(\widetilde x_A|_U)\cap P(U))|\cdot |\pi_0(Z_{G(W)}(y))|.
\] Since \[
|\pi_0(Z_{G(U)}(\widetilde x_A|_U))|=4
\qquad\text{and}\qquad
|\pi_0(Z_{G(U)}(\widetilde x_A|_U)\cap P(U))|=2,
\] the common \(W\)-factor cancels and yields \[
m(G\cdot \widetilde x_A,P)=2.
\] Because \(\widetilde x_A\) is \(P\)-conjugate to \(x_A\), this is
exactly \[
m(G\cdot x_A,P)=2.
\]

This completes the case \(\epsilon=-1\) and \(d\) odd, and hence the
proof.

\section{theorem thm:main}\label{theorem-thmmain}

\subsection{statement}\label{statement-8}

\section{Short-Kappa Branch of Maximal-Parabolic Induction - Computation
Problem}\label{short-kappa-branch-of-maximal-parabolic-induction---computation-problem}

Let \[
F=\mathbb R.
\] Let \[
\epsilon\in\{+1,-1\}.
\]

\subsection{Setup}\label{setup}

Let \((V_0,\langle\cdot,\cdot\rangle_0)\) be a finite-dimensional
\(F\)-vector space with a non-degenerate bilinear form satisfying \[
\langle v,w\rangle_0=\epsilon\langle w,v\rangle_0
\qquad \forall v,w\in V_0.
\] Set \[
G_0:=\{g\in \mathrm{GL}(V_0):\langle gv,gw\rangle_0=\langle v,w\rangle_0\},
\] and \[
\mathfrak g_0:=\{Y\in\mathrm{End}_F(V_0):
\langle Yv,w\rangle_0+\langle v,Yw\rangle_0=0\ \forall v,w\in V_0\}.
\]

Let \(X_0\in\mathfrak g_0\) be nilpotent, and let \[
\mathcal O_0:=G_0\cdot X_0.
\] Put \[
c_1:=\dim_F\ker X_0,
\qquad
c_2:=\dim_F(\ker X_0\cap\operatorname{Im}X_0).
\] Equivalently, \[
c_2=\dim_F(\ker X_0^2/\ker X_0).
\]

Fix an integer \(\kappa\), and let \(E,E'\) be \(\kappa\)-dimensional
\(F\)-vector spaces with a non-degenerate pairing \[
\lambda:E\times E'\to F.
\] Assume \[
\dim_F\ker X_0>\dim_F E=\kappa\ge \dim_F(\ker X_0^2/\ker X_0).
\] Equivalently, \[
c_2\le \kappa<c_1.
\] Put \[
d:=c_1-\kappa.
\] Set \[
V:=E\oplus V_0\oplus E',
\] and define a bilinear form on \(V\) by \[
\langle e+v+e',\ f+w+f'\rangle
:=\lambda(e,f')+\epsilon\lambda(f,e')+\langle v,w\rangle_0.
\] Let \[
G:=\{g\in\mathrm{GL}(V):\langle gx,gy\rangle=\langle x,y\rangle\},
\] and \[
\mathfrak g:=\{Y\in\mathrm{End}_F(V):
\langle Yx,y\rangle+\langle x,Yy\rangle=0\ \forall x,y\in V\}.
\]

Let \[
P:=\{g\in G:g(E)=E\}.
\] Since \(E^\perp=E\oplus V_0\), define \[
\mathfrak u:=\{Y\in\mathfrak g:
Y(E)=0,\quad
Y(E^\perp)\subseteq E,\quad
Y(V)\subseteq E^\perp
\}.
\] Regard \(\mathcal O_0\) as a subset of \(\mathfrak g\) by letting
each operator in \(\mathcal O_0\) act on the \(V_0\)-summand and act as
\(0\) on \(E\oplus E'\).

Define the induced set \[
\operatorname{Ind}_P^G(\mathcal O_0)
:=\overline{G\cdot(\mathcal O_0+\mathfrak u)}\subseteq\mathfrak g,
\] where the closure is taken in the usual real topology on the
finite-dimensional real vector space \(\mathfrak g\).

A \(G\)-orbit
\(\mathcal O\subseteq\operatorname{Ind}_P^G(\mathcal O_0)\) is called
induced from \(\mathcal O_0\) via \(P\) if it is maximal, among the
\(G\)-orbits contained in \(\operatorname{Ind}_P^G(\mathcal O_0)\), for
the closure order \[
\mathcal O_1\preceq \mathcal O_2
\iff
\mathcal O_1\subseteq\overline{\mathcal O_2}.
\]

For \(x\in\mathfrak g\), write \[
Z_G(x):=\{g\in G:gxg^{-1}=x\}.
\] For a subgroup \(H\) of \(G\), let \(H^\circ_{\mathrm{top}}\) denote
its identity component for the usual real topology, and set \[
\pi_0(H):=H/H^\circ_{\mathrm{top}}.
\] For \(x\in\mathcal O_0+\mathfrak u\), define the multiplicity by \[
m(G\cdot x,P):=
\frac{|\pi_0(Z_G(x))|}{|\pi_0(Z_G(x)\cap P)|}.
\]

\subsection{The kernel quotient}\label{the-kernel-quotient}

Since \(X_0\) is skew-adjoint, one has \[
(\operatorname{Im}X_0)^\perp=\ker X_0.
\] Therefore the restriction of \(\langle\cdot,\cdot\rangle_0\) to
\(\ker X_0\) has radical \[
\ker X_0\cap\operatorname{Im}X_0.
\] Define the kernel quotient \[
K_{X_0}:=\ker X_0/(\ker X_0\cap\operatorname{Im}X_0).
\] The restriction of \(\langle\cdot,\cdot\rangle_0\) to \(\ker X_0\)
descends to a bilinear form on \(K_{X_0}\), denoted by \[
\langle\cdot,\cdot\rangle_K.
\] Its dimension is \[
\dim_F K_{X_0}=c_1-c_2.
\] The assumption \(c_2\le\kappa<c_1\) is equivalent to \[
0<d\le \dim_F K_{X_0}.
\]

Let \(\mathscr A\) be the set of \(d\)-dimensional subspaces \[
A\subset K_{X_0}.
\] For \(A\in\mathscr A\), write \(\operatorname{rad}(A)\) for the
radical of the restricted form \(\langle\cdot,\cdot\rangle_K|_A\). Let
\(\mathscr A^{\mathrm{mr}}\subseteq\mathscr A\) be the maximal-rank
locus: this is the set of all \(A\) for which the restricted form on
\(A\) has the largest possible rank among \(d\)-dimensional subspaces of
\(K_{X_0}\).

Two elements \(A_1,A_2\in\mathscr A^{\mathrm{mr}}\) are called isometric
inside \(K_{X_0}\) if there is an isometry \[
h\in\operatorname{Isom}(K_{X_0},\langle\cdot,\cdot\rangle_K)
\] such that \[
h(A_1)=A_2.
\]

\subsection{\texorpdfstring{The slice
\(X_0+\mathfrak u\)}{The slice X\_0+\textbackslash mathfrak u}}\label{the-slice-x_0mathfrak-u}

With respect to \[
V=E\oplus V_0\oplus E',
\] every element of \(X_0+\mathfrak u\) has the form \[
X_{C,B}:=
\begin{bmatrix}
0&C^\vee&B\\
0&X_0&C\\
0&0&0
\end{bmatrix}.
\] Here \[
C:E'\to V_0,
\qquad
B:E'\to E,
\] and \(C^\vee:V_0\to E\) is determined by \[
\lambda(C^\vee v,a')=-\langle v,Ca'\rangle_0
\qquad \forall v\in V_0,\ a'\in E'.
\] The condition on \(B\) is \[
\lambda(Ba',b')+\epsilon\lambda(Bb',a')=0
\qquad \forall a',b'\in E'.
\]

\subsection{Problem}\label{problem}

Compute the maximal \(G\)-orbits in \[
\operatorname{Ind}_P^G(\mathcal O_0)
\] and compute their multiplicities with respect to the maximal
parabolic subgroup \(P\).

The final answer should give a parametrization of the maximal induced
orbits in terms of the isometry classes of the \(d\)-dimensional forms
obtained from subspaces of \[
K_{X_0}=\ker X_0/(\ker X_0\cap\operatorname{Im}X_0).
\] For each such orbit \(\mathcal O\), compute \[
m(\mathcal O,P)
=
\frac{|\pi_0(Z_G(x))|}
{|\pi_0(Z_G(x)\cap P)|},
\] where \(x\in\mathcal O\cap(\mathcal O_0+\mathfrak u)\) is any
representative used in the calculation.

Give a self-contained proof from the maximal parabolic action on the
slice \(X_0+\mathfrak u\). Do not use any diagram parametrization.

\subsection{proof}\label{proof-8}

Lemmas \texttt{lem:x0-geometry-and-k},
\texttt{lem:parametrize-x0-plus-u}, and
\texttt{lem:unipotent-conjugation} identify the slice and isolate the
intrinsic kernel-quotient datum carried by the surjective map \[
\phi_C=C^\vee|_{\ker X_0}:\ker X_0\to E.
\] Proposition \texttt{prop:generic-normal-form} shows that on the open
dense locus where \(\phi_C\) is surjective and its kernel is transverse
to \(R=\ker X_0\cap \operatorname{Im} X_0\), every slice point is
\(P\)-conjugate to a normal form \(x_A\) indexed by a \(d\)-dimensional
subspace \[
A\subset K_{X_0}\simeq V_1.
\]

Among these subspaces, the open dense maximal-rank locus \[
\mathscr A^{\mathrm{mr}}
\] is exactly the one that survives as the maximal part of the induced
set. Lemma \texttt{lem:maximal-rank-classes} shows that only finitely
many isometry classes occur there. Proposition
\texttt{prop:induced-orbits} proves that the corresponding \(G\)-orbits
\[
G\cdot x_A \qquad (A\in \mathscr A^{\mathrm{mr}})
\] are precisely the maximal \(G\)-orbits in \[
\operatorname{Ind}_P^G(\mathcal O_0),
\] and that two such maximal orbits coincide exactly when the
corresponding subspaces are isometric inside
\((K_{X_0},\langle\cdot,\cdot\rangle_K)\).

Finally, Proposition \texttt{prop:multiplicity} computes the
component-group quotient for every maximal representative and gives \[
m(G\cdot x_A,P)=
\begin{cases}
1,& \epsilon=+1,\\[0.4ex]
1,& \epsilon=-1\ \text{and}\ d\ \text{even},\\[0.4ex]
2,& \epsilon=-1\ \text{and}\ d\ \text{odd}.
\end{cases}
\]

Therefore:

\begin{enumerate}
\def\labelenumi{\arabic{enumi}.}
\item
  The maximal \(G\)-orbits in \(\operatorname{Ind}_P^G(\mathcal O_0)\)
  are parametrized by the isometry classes inside \(K_{X_0}\) of the
  maximal-rank \(d\)-dimensional subspaces \(A\subset K_{X_0}\).
\item
  If \(A\) is such a subspace and \(x_A\in \mathcal O_0+\mathfrak u\) is
  the normal representative constructed above, then \[
  \mathcal O_A:=G\cdot x_A
  \] is maximal in the induced set, and every maximal induced orbit is
  obtained in this way.
\item
  For every maximal induced orbit, \[
  m(\mathcal O_A,P)=
  \begin{cases}
  1,& \epsilon=+1,\\[0.4ex]
  1,& \epsilon=-1\ \text{and}\ d\ \text{even},\\[0.4ex]
  2,& \epsilon=-1\ \text{and}\ d\ \text{odd}.
  \end{cases}
  \]
\end{enumerate}

This proves the problem.

\end{document}
